% Compilar en Overleaf con pdfLaTeX
\documentclass[aspectratio=169,10pt]{beamer}
\usetheme{Madrid}
\usecolortheme{whale}
\usepackage[utf8]{inputenc}
\usepackage[spanish]{babel}
\usepackage{listings}
\usepackage{xcolor}

\lstdefinestyle{codigo}{
  basicstyle=\ttfamily\footnotesize,
  breaklines=true,
  frame=single,
  backgroundcolor=\color{gray!10},
  keepspaces=true,
  showstringspaces=false
}

\lstset{
  basicstyle=\ttfamily\footnotesize,
  breaklines=true,
  frame=single,
  backgroundcolor=\color{gray!10},
  keywordstyle=\color{blue},
  commentstyle=\color{green!50!black}
}

\title{Contenedores con Docker}
\subtitle{Dockerfile, Docker Compose y API REST en Laravel}
\author{Albeiro Montes}
\institute{Ingeniería de Software II / Proyecto CRUD Laravel}
\date{\today}

\begin{document}

\begin{frame}
  \titlepage
\end{frame}

\begin{frame}{Agenda}
  \tableofcontents
\end{frame}

% ============================================================
\section{¿Por qué contenedores?}
% ============================================================

\begin{frame}{Del código a producción}
  \begin{itemize}
    \item En local funciona; en el servidor de un compañero, no.
    \item Dependencias distintas: versión de PHP, extensiones, MySQL, rutas.
    \item Objetivo: \textbf{mismo entorno} para desarrollo, pruebas y despliegue.
  \end{itemize}
  \vspace{0.5em}
  \begin{block}{Idea clave}
    Empaquetar la aplicación con todo lo que necesita para ejecutarse, de forma aislada y reproducible.
  \end{block}
\end{frame}

\begin{frame}{Máquina virtual vs contenedor}
  \begin{columns}[T]
    \column{0.48\textwidth}
    \textbf{VM}
    \begin{itemize}
      \item SO completo invitado
      \item Hypervisor
      \item Arranque lento, más RAM
    \end{itemize}
    \column{0.48\textwidth}
    \textbf{Contenedor}
    \begin{itemize}
      \item Comparte el kernel del host
      \item Solo procesos y filesystem aislados
      \item Ligero, segundos para levantar
    \end{itemize}
  \end{columns}
  \vspace{0.8em}
  \centering
  \small Host $\rightarrow$ Docker Engine $\rightarrow$ Contenedores (Apache + PHP, MySQL)
\end{frame}

\begin{frame}{Vocabulario Docker}
  \begin{description}
    \item[Imagen] Plantilla de solo lectura (capas). Ej.: \texttt{mysql:8.0}, \texttt{php:8.2-apache}.
    \item[Contenedor] Instancia en ejecución de una imagen.
    \item[Dockerfile] Receta para \textbf{construir} una imagen propia (tu app PHP + extensiones).
    \item[Docker Compose] Archivo YAML que define \textbf{varios servicios} y cómo se conectan.
    \item[Volumen] Carpeta del host montada en el contenedor (el código de Laravel).
    \item[Red] Permite que \texttt{app} resuelva el host \texttt{db} por nombre.
  \end{description}
\end{frame}

% ============================================================
\section{Arquitectura del proyecto}
% ============================================================

\begin{frame}{API Laravel en contenedores}
  \textbf{Caso: CRUD / API REST} con Laravel, Apache y MySQL.
  \vspace{0.5em}
  \begin{itemize}
    \item \textbf{Apache + PHP}: un solo servicio (\texttt{app}) recibe HTTP en el puerto 80 y ejecuta Laravel.
    \item El documento raíz es \texttt{public/}; \texttt{mod\_rewrite} y \texttt{.htaccess} envían las rutas a \texttt{index.php}.
    \item \textbf{MySQL}: base de datos del CRUD.
  \end{itemize}
  \vspace{0.5em}
  \begin{alertblock}{Flujo de una petición}
    Cliente $\rightarrow$ Apache $\rightarrow$ Laravel $\rightarrow$ MySQL
  \end{alertblock}
\end{frame}

\begin{frame}[fragile]{Imágenes base a descargar (pull)}
  Antes de construir, se descargan las imágenes que usa el proyecto:
\begin{lstlisting}[language=bash]
docker pull php:8.2-apache
docker pull composer:2
docker pull mysql:8.0
\end{lstlisting}
  \vspace{0.3em}
  \begin{itemize}
    \item \texttt{php:8.2-apache}: PHP 8.2 con Apache (base del Dockerfile).
    \item \texttt{composer:2}: aporta el binario de Composer (build); \texttt{vendor/} se genera al arrancar el contenedor.
    \item \texttt{mysql:8.0}: base de datos, usada tal cual por Compose.
  \end{itemize}
  \vspace{0.2em}
  \small La imagen de \textbf{tu aplicación} no se descarga: se genera con \texttt{docker compose build}.
\end{frame}

% ============================================================
\section{Paso a paso del proyecto}
% ============================================================

\begin{frame}[fragile]{Paso 1: Requisitos en el equipo}
  \begin{enumerate}
    \item Instalar \textbf{Docker Desktop} (Windows/Mac) o Docker Engine + Compose (Linux).
    \item Verificar:
  \end{enumerate}
\begin{lstlisting}[language=bash]
docker --version
docker compose version
\end{lstlisting}
  \begin{enumerate}
    \setcounter{enumi}{2}
    \item Tener el \textbf{proyecto Laravel (API) ya creado en el host} (Composer/PHP local).
  \end{enumerate}
\end{frame}

\begin{frame}[fragile]{Paso 2: Laravel en el host + preparar Docker}
  \textbf{Partimos de una app existente} en la carpeta del repo (CRUD / \texttt{routes/api.php}).
  \vspace{0.4em}
  \begin{enumerate}
    \item Añadir \texttt{dockerfile}, \texttt{docker-compose.yml} y \texttt{entrypoint.sh} en la raíz del proyecto Laravel.
    \item Crear \texttt{.env} para cuando la app corre \textbf{dentro} de Compose (\texttt{DB\_HOST=db}, usuario \texttt{root}):
  \end{enumerate}
\begin{lstlisting}
APP_NAME=Laravel
APP_ENV=local
APP_KEY=
APP_DEBUG=true
APP_URL=http://localhost:8080

DB_CONNECTION=mysql
DB_HOST=db
DB_PORT=3306
DB_DATABASE=laravel
DB_USERNAME=root
DB_PASSWORD=secret
\end{lstlisting}
\end{frame}

\begin{frame}{Paso 3: Dockerfile (visión general)}
  El Dockerfile \textbf{parte de PHP con Apache} e instala lo que Laravel necesita:
  \begin{itemize}
    \item Extensiones: \texttt{pdo\_mysql}, \texttt{mbstring}, \texttt{bcmath}, \texttt{gd}.
    \item \texttt{mod\_rewrite} y \texttt{AllowOverride All} para que funcione \texttt{public/.htaccess}.
    \item Document root de Apache en \texttt{/var/www/html/public}.
    \item Composer (imagen \texttt{composer:2}) instala \texttt{vendor/} al \textbf{arrancar} el contenedor, porque el volumen del host reemplaza el código de la imagen.
    \item Permisos de \texttt{www-data} en \texttt{storage} y \texttt{bootstrap/cache}.
  \end{itemize}
  \vspace{0.3em}
  Cada instrucción \texttt{RUN} añade una capa a la imagen.
\end{frame}

\begin{frame}[fragile]{Dockerfile de referencia}
\begin{lstlisting}[language=bash,basicstyle=\ttfamily\scriptsize]
FROM php:8.2-apache

RUN apt-get update && apt-get install -y \
    git unzip libpng-dev libonig-dev libxml2-dev \
    && docker-php-ext-install pdo_mysql mbstring bcmath gd \
    && a2enmod rewrite \
    && sed -ri -e 's!/var/www/html!/var/www/html/public!g' \
       /etc/apache2/sites-available/*.conf \
    && sed -ri -e 's!AllowOverride None!AllowOverride All!g' \
       /etc/apache2/apache2.conf \
    && rm -rf /var/lib/apt/lists/*

COPY --from=composer:2 /usr/bin/composer /usr/bin/composer
COPY entrypoint.sh /usr/local/bin/laravel-entrypoint.sh
RUN chmod +x /usr/local/bin/laravel-entrypoint.sh

WORKDIR /var/www/html
EXPOSE 80
ENTRYPOINT ["laravel-entrypoint.sh"]
\end{lstlisting}
  \small \texttt{ENTRYPOINT} ejecuta \texttt{entrypoint.sh} cada vez que arranca el contenedor \texttt{app}.
\end{frame}

\begin{frame}{entrypoint.sh: ¿por qué existe?}
  En \texttt{docker-compose.yml} montamos \texttt{.:/var/www/html}:
  \begin{itemize}
    \item El código del host \textbf{sustituye} lo copiado en la imagen al construirla.
    \item Si en el PC no hay carpeta \texttt{vendor/}, Laravel falla con \texttt{autoload.php}.
    \item Solución: al \textbf{iniciar} el contenedor, un script ejecuta \texttt{composer install} y luego Apache.
  \end{itemize}
  \vspace{0.4em}
  \begin{block}{Montaje en el Dockerfile}
    \texttt{COPY entrypoint.sh} $\rightarrow$ \texttt{/usr/local/bin/laravel-entrypoint.sh} \\
    \texttt{ENTRYPOINT ["laravel-entrypoint.sh"]}
  \end{block}
  \vspace{0.3em}
  \small En Overleaf, el listado usa estilo \texttt{codigo} (sin resaltado \texttt{bash}) para que el shebang \texttt{\#!/bin/bash} no desaparezca.
\end{frame}

\begin{frame}[fragile]{entrypoint.sh: código completo}
\begin{lstlisting}[style=codigo]
#!/bin/bash
set -e

cd /var/www/html

composer install --no-interaction --prefer-dist --optimize-autoloader

chown -R www-data:www-data storage bootstrap/cache

exec docker-php-entrypoint apache2-foreground
\end{lstlisting}
\end{frame}

\begin{frame}{entrypoint.sh: línea por línea}
  \small
  \begin{description}
    \item[\texttt{\#!/bin/bash}] Intérprete del script.
    \item[\texttt{set -e}] Si un comando falla, el contenedor se detiene (no arranca Apache a medias).
    \item[\texttt{cd /var/www/html}] Directorio de Laravel (coincide con el volumen montado).
    \item[\texttt{composer install ...}] Crea o actualiza \texttt{vendor/} \textbf{dentro del contenedor}.
    \item[\texttt{chown ...}] Laravel puede escribir en \texttt{storage} y \texttt{bootstrap/cache}.
    \item[\texttt{exec docker-php-entrypoint ...}] Arranca Apache en primer plano (proceso principal del contenedor).
  \end{description}
\end{frame}

\begin{frame}{Paso 4: Docker Compose (visión general)}
  \texttt{docker-compose.yml} declara dos servicios:
  \begin{itemize}
    \item \texttt{app}: \texttt{build: .} usa el Dockerfile. Publica el puerto \texttt{8080} del host en el \texttt{80} de Apache.
    \item \texttt{db}: imagen \texttt{mysql:8.0}, base \texttt{laravel} y contraseña de root \texttt{secret}.
  \end{itemize}
  \vspace{0.4em}
  El volumen \texttt{.:/var/www/html} monta el código del host dentro del contenedor. Ambos servicios comparten red, así Laravel llega a MySQL con \texttt{DB\_HOST=db}.
\end{frame}

\begin{frame}[fragile]{docker-compose.yml}
\begin{lstlisting}[language=bash]
services:
  app:
    build: .
    ports:
      - "8080:80"
    volumes:
      - .:/var/www/html
    depends_on:
      - db

  db:
    image: mysql:8.0
    environment:
      MYSQL_DATABASE: laravel
      MYSQL_ROOT_PASSWORD: secret
\end{lstlisting}
\end{frame}

\begin{frame}[fragile]{Paso 5: Levantar el entorno}
\begin{lstlisting}[language=bash]
# Descargar imagenes y construir la app
docker compose pull
docker compose build

# Al arrancar, el contenedor ejecuta composer install
docker compose up

# Ver estado
docker compose ps
\end{lstlisting}
  Abrir en el navegador: \texttt{http://localhost:8080}
\end{frame}

\begin{frame}[fragile]{Paso 6: Laravel dentro del contenedor}
\begin{lstlisting}[language=bash]
# Clave de aplicacion (escribe APP_KEY en .env)
docker compose exec app php artisan key:generate

# Migraciones (tablas del CRUD)
docker compose exec app php artisan migrate

# Rutas API (ejemplo)
docker compose exec app php artisan route:list
\end{lstlisting}
  Los endpoints REST viven en \texttt{routes/api.php} y los controladores en \texttt{app/Http/Controllers}.
\end{frame}

\begin{frame}[fragile]{Paso 7: Desarrollo diario}
\begin{lstlisting}[language=bash]
# Logs de Apache y de PHP
docker compose logs -f app

# Entrar al contenedor de la aplicacion
docker compose exec app bash

# Parar los contenedores
docker compose stop

# Parar y eliminar contenedores
docker compose down
\end{lstlisting}
  \vspace{0.3em}
  \small MySQL no tiene un volumen con nombre: \texttt{docker compose down} elimina también los datos de la base.
\end{frame}

% ============================================================
\section{Comandos útiles}
% ============================================================

\begin{frame}[fragile]{Imágenes y contenedores}
\begin{lstlisting}[language=bash]
docker pull php:8.2-apache
docker pull composer:2
docker pull mysql:8.0

docker images
docker ps
docker ps -a
\end{lstlisting}
\end{frame}

\begin{frame}[fragile]{Depuración rápida}
  \begin{itemize}
    \item \textbf{La página no abre}: el contenedor \texttt{app} no está en marcha o el puerto publicado no es \texttt{8080}.
    \item \textbf{vendor/autoload.php}: falta \texttt{vendor/}; reconstruir y levantar (\texttt{entrypoint.sh} ejecuta \texttt{composer install}) o \texttt{docker compose exec app composer install}.
    \item \textbf{Error 500}: falta \texttt{.env} o \texttt{APP\_KEY}. Generar la clave con \texttt{artisan key:generate}.
    \item \textbf{SQLSTATE connection refused}: \texttt{DB\_HOST} debe ser \texttt{db}, el usuario \texttt{root} y la contraseña \texttt{secret}. MySQL puede tardar unos segundos en aceptar conexiones.
    \item \textbf{Permisos Laravel}: \texttt{storage} y \texttt{bootstrap/cache} deben ser escribibles por \texttt{www-data}.
  \end{itemize}
\begin{lstlisting}[language=bash]
docker compose exec db mysql -uroot -psecret laravel
\end{lstlisting}
\end{frame}

% ============================================================
\section{Resumen}
% ============================================================

\begin{frame}{Resumen}
  \begin{itemize}
    \item Contenedor = entorno reproducible; Compose = orquestación local de varios servicios.
    \item \textbf{Pull}: \texttt{php:8.2-apache}, \texttt{composer:2} y \texttt{mysql:8.0}. La app se construye con el Dockerfile.
    \item Apache sirve \texttt{public/} y Laravel se conecta con \texttt{DB\_HOST=db}, usuario \texttt{root} y contraseña \texttt{secret}.
    \item Flujo: \texttt{dockerfile} + \texttt{entrypoint.sh} + Compose $\rightarrow$ \texttt{up} (Composer en el contenedor) $\rightarrow$ \texttt{key:generate} y \texttt{migrate} $\rightarrow$ API en el puerto 8080.
  \end{itemize}
\end{frame}

\begin{frame}{Referencias}
  \small
  \begin{itemize}
    \item Documentación Docker: \url{https://docs.docker.com/}
    \item Imágenes oficiales PHP, Composer y MySQL en Docker Hub.
    \item Laravel: despliegue y configuración \texttt{.env}.
  \end{itemize}
  \vspace{1em}
  \centering
  ¿Preguntas?
\end{frame}

\end{document}
